\labsection{4}{CPU scheduling: compare policies on one workload}{sec:lab04}
\begin{labproject}{Make scheduling policy visible through order, response, and waiting time}

\textbf{Coverage.} Chapter~\ref{ch:scheduling}: scheduling metrics and policy rules.\\
\textbf{Goal.} Run four policies on the same bursts so only the policy changes.\\
\textbf{Explain.} Use the trace to explain why order or time quantum changes a metric.

\section*{Source code}
\labsourcecard{\texttt{fcfs.c}}{Run processes in the supplied order.}{Waiting time is the accumulated burst time of earlier processes.}
\lstinputlisting[style=utpcnumbered]{all_experiments/Lab04/fcfs.c}

\labsourcecard{\texttt{sjf.c}}{Sort the same workload by shortest burst first.}{The sorting step changes execution order; the metric loop is otherwise very similar to FCFS.}
\lstinputlisting[style=utpcnumbered]{all_experiments/Lab04/sjf.c}

\labsourcecard{\texttt{priority.c}}{Sort by priority, where a smaller number means higher priority.}{The code uses PID only as a tie-breaker so equal priorities are deterministic.}
\lstinputlisting[style=utpcnumbered]{all_experiments/Lab04/priority.c}

\labsourcecard{\texttt{round\_robin.c}}{Give each unfinished process a bounded time slice.}{\texttt{remaining[]} tracks work left; \texttt{first\_run[]} gives response time when all arrivals are 0.}
\lstinputlisting[style=utpcnumbered]{all_experiments/Lab04/round_robin.c}

\section*{Task 1 --- Compile all policies}
\begin{lstlisting}[style=utpbash]
cd all_experiments/Lab04
for f in fcfs sjf priority round_robin; do
  gcc -Wall -Wextra -std=c11 "$f.c" -o "$f"
done
\end{lstlisting}

\section*{Task 2 --- Baseline comparison}
\begin{lstlisting}[style=utpbash]
./fcfs
./sjf
./priority
./round_robin 2
\end{lstlisting}
Complete a table with execution order, first process to finish, average waiting time when printed, and one observation about responsiveness.

\begin{utptable}
\renewcommand{\arraystretch}{1.18}
\begin{tabularx}{\linewidth}{@{}>{\bfseries\leavevmode\color{UTPNavy}}m{0.18\linewidth}
  *{3}{>{\raggedright\arraybackslash}X}@{}}
\toprule
\rowcolor{UTPPaleBlue!92}
\utpth{Policy} & \utpth{First execution order} & \utpth{Waiting/response evidence} & \utpth{What the rule favours} \\
\midrule
FCFS & & & \\
SJF & & & \\
Priority & & & \\
RR, $q=2$ & & & \\
\bottomrule
\end{tabularx}
\end{utptable}

\section*{Task 3 --- Verify one result by hand}
For FCFS, draw a one-line Gantt chart for bursts 8, 4, 2, 1. Compute the waiting time of each process before comparing your arithmetic with the program.

\section*{Task 4 --- Change only the Round Robin quantum}
\begin{lstlisting}[style=utpbash]
./round_robin 1
./round_robin 2
./round_robin 4
\end{lstlisting}
Record the first response time of each process for each quantum. Explain the trade-off visible in the Gantt output: smaller turns provide earlier rotation but create more slices.

\begin{quickcheck}
\begin{enumerate}
  \item Which scheduling metric measures time spent waiting in the ready queue?
  \item Why can Round Robin improve response time while increasing the number of context switches?
  \item Why must all policies be compared using the same workload?
\end{enumerate}
\end{quickcheck}

\begin{verification}
\begin{itemize}
  \item All four policies run on the same bursts, so only the policy differs.
  \item The FCFS Gantt chart and its waiting times are checked by hand against the program.
  \item The Round Robin first-response times are recorded for $q=1,2,4$.
  \item Each explanation names the policy rule first and the metric consequence second.
\end{itemize}
\end{verification}

\begin{labdeliverable}
Submit the four-policy comparison table, your hand-checked FCFS Gantt chart, and a small $q=1,2,4$ response-time table. Your explanation should identify a policy rule first and a metric consequence second.
\end{labdeliverable}
\end{labproject}
